\section{Afinidad con el hardware: el algoritmo debe poder paralelizarse, usar multiplicación de matrices y escalar}

Después de la arqueología de algoritmos, esta sección aborda el último umbral para que una idea llegue a la práctica: el hardware. Si una idea no puede paralelizarse, no puede escribirse como un kernel eficiente ni aprovechar la capacidad de multiplicación de matrices de las GPU modernas, difícilmente llegará a los modelos frontier.

La última línea técnica es la afinidad con el hardware. Yang Songlin (杨松琳) sostiene que un algoritmo no solo debe tener sentido desde el punto de vista de machine learning, sino que también debe poder paralelizarse, escalar e implementarse de forma eficiente en el hardware. Transformer tuvo éxito en su momento no solo por su gran capacidad expresiva, sino también porque tanto attention como FFN contienen abundantes multiplicaciones de matrices, lo que les da una mayor afinidad con el hardware que RNN/LSTM.

\begin{figure}[H]
\centering
\includegraphics[width=0.96\textwidth]{hardware-friendly-algorithm.png}
\caption{Algoritmos afines al hardware: además de tener sentido matemático, deben poder paralelizarse, usar multiplicación de matrices y escalar. Diagrama conceptual propio, elaborado a partir de la conversación de 01:29:50--01:42:23.}
\end{figure}

\begin{knowledgebox}{Lectura de la figura: la investigación actual en arquitecturas es un diseño conjunto de algoritmos y sistemas}
Que sea matemáticamente sólido es solo el primer paso; también se necesita un algoritmo paralelo, como chunkwise; debe poder aprovechar la multiplicación de matrices; debe tener una implementación como kernel, y, por último, debe escalar al tamaño de los modelos grandes. Si falta cualquiera de estos eslabones, al algoritmo le será difícil llegar a los modelos frontier.
\end{knowledgebox}

\subsection{Chunkwise Algorithm y paralelismo}

Esta subsección concreta la afinidad con el hardware en forma de algoritmos paralelos. La atención lineal suele tener una forma recurrente, y la recurrencia por sí misma no se presta al paralelismo masivo en GPU, por lo que es necesario reescribirla en una forma que pueda paralelizarse por bloques.

Linear Attention suele tener una forma recurrente, pero la recurrencia es intrínsecamente difícil de paralelizar. Para entrenarla de manera eficiente en las GPU modernas se requieren algoritmos paralelos como el chunkwise algorithm, que divide las secuencias largas en bloques: conserva la estructura recurrente y, al mismo tiempo, aprovecha el cómputo en paralelo. Su papel es similar al de FlashAttention respecto a Softmax Attention: no cambia el objetivo del modelo, sino que hace que el algoritmo pueda ejecutarse en el hardware.

\begin{importantbox}{La importancia de los algoritmos paralelos}
Sin un algoritmo paralelo eficiente, incluso una variante de attention muy elegante en teoría no puede llegar al entrenamiento a gran escala. La innovación en arquitecturas y la optimización de sistemas no son etapas sucesivas: juntas determinan si una idea puede llevarse a la práctica.
\end{importantbox}

\subsection{La coordinación con el hardware en DeepSeek}

A continuación se toma DeepSeek como ejemplo de coordinación con el hardware. Su valor no está solo en haber elegido sparse, sino en que la forma de implementar sparse se ajusta a FP8, a la multiplicación de matrices y a una ruta de hardware sin softmax.

En la entrevista se considera que DeepSeek da mucha importancia a la coordinación entre hardware y algoritmos. El indexer de DeepSeek Sparse Attention puede calcular el attention score / logit en FP8, eliminar las costosas operaciones softmax/exp y dejar principalmente multiplicaciones de matrices y la selección Top-K. Esto lo acerca más a los principios del hardware y facilita la coordinación con el equipo de infra.

\begin{figure}[H]
\centering
\includegraphics[width=0.94\textwidth]{deepseek-hardware-codesign.png}
\caption{La coordinación con el hardware en DeepSeek: la selección dispersa, FP8 y la multiplicación de matrices sirven en conjunto a la eficiencia. Diagrama conceptual propio, elaborado a partir de la conversación de 01:29:50--01:42:23.}
\end{figure}

\begin{knowledgebox}{Lectura de la figura: la coordinación con el hardware no consiste solo en escribir kernels}
Del lado del algoritmo, Sparse / Indexer / Top-K reducen el cómputo; del lado del hardware, la eficiencia aumenta con FP8, la multiplicación de matrices y la eliminación de las costosas operaciones softmax/exp. El verdadero co-design considera al mismo tiempo los objetivos del algoritmo y los principios del hardware.
\end{knowledgebox}

\subsection{Resumen de la sección}

La investigación futura en arquitecturas deberá cumplir tres condiciones a la vez: ser sólida desde el punto de vista del aprendizaje automático, ser explicable matemáticamente y, en el hardware, poder paralelizarse y usar multiplicación de matrices. De lo contrario, por elegante que sea un algoritmo, difícilmente llegará a los modelos frontier.
